\documentclass[../../../include/open-logic-section]{subfiles}

\begin{document}

\olfileid{cmp}{thy}{int}
\olsection{引言}

逻辑学中称为\emph{可计算性理论}的分支，研究函数与集合的可计算性或相对可计算性问题。由于 Kleene的影响，该学科曾被称为\emph{递归理论}，如今这两个名称都常被使用。% Part: computability
% Chapter: computability-theory
% Section: introduction

若函数~$f\colon \Nat \pto \Nat$ 能在某种计算模型中计算，我们就称其为\emph{部分可计算的}。若~$f$ 是全函数，则简称~$f$ 是\emph{可计算的}。若关系~$R$ 的特征函数~$\Char{R}$ 是可计算的，则称该关系也是可计算的。若 $f$ 和~$g$ 是偏函数，我们用 $f(x) \fdefined$ 表示 $f$ 在~$x$ 处有定义，即 $x$ 属于~$f$ 的定义域；反之，用 $f(x) \fundefined$ 表示 $f$ 在~$x$ 处无定义。我们用 $f(x) \simeq g(x)$ 表示：要么 $f(x)$ 和 $g(x)$ 均无定义，要么均有定义且相等。

我们可以不依赖特定的计算模型来研究可计算性理论。为此，只需证明存在一个通用部分可计算函数~$\fn{Un}(k, x)$。由此我们可以枚举部分可计算函数。我们将采用记号~$\cfind{k}$ 表示第~$k$ 个一元部分可计算函数，其定义为$\cfind{k}(x) \simeq \fn{Un}(k, x)$。（Kleene 曾为此使用 $\{ k \}$，但该记号近来已不多用。）稍作推广，我们可以一致地枚举任意元数的部分可计算函数，并用 $\cfind{k}[n]$ 表示第~$k$ 个 $n$ 元部分递归函数。

如果 $f(\vec x, y)$ 是全函数或偏函数，那么$\umin{y}{f
(\vec x, y)}$ 是关于自变元~$\vec x$ 的函数，它返回使 $f(\vec x, y) = 0$ 成立的最小~$y$，且假定 $f(\vec x, 0)$, \dots, $f(\vec x, y-1)$ 均有定义；若不存在这样的 $y$，则$\umin{y}{f (\vec x, y)}$ 无定义。如果 $R(\vec x, y)$ 是关系，则$\umin{y}{R(\vec x, y)}$ 定义为使 $R(\vec x, y)$ 为真的最小~$y$；换言之，即满足$1 \tsub \Char{R}(\vec x, y) = 0$的最小~$y$。

要证明一个函数是可计算的，有两条途径：
\begin{enumerate}
\item \emph{严格地}：明确地描述一个Turing机或部分递归函数，并证明它计算了目标函数；
\item \emph{非形式地}：描述一个计算它的算法，并诉诸 Church论题。
\end{enumerate}

两者之间并无严格的界限；对算法的详细描述应提供足够的信息，使人原则上能相对清楚地看出如何设计出正确的Turing机或部分递归定义序列。完全严格的定义往往缺乏启发性，我们将尝试在这两种方法之间寻找合理的折中；简言之，我们将努力提供既直观又严格的证明，以确保相应的精确定义能够得出。

\end{document}